% Neso videos on modular exponentiation, part 1 (base reduction) and part 2 (repeated squaring) (notes8.md); Forouzan 9.4
% (square-and-multiply, right-to-left); Stallings (fast modular exponentiation, left-to-right). Every number recomputed by script (scratchpad/r4/verify.py, v2.py).
\section{Modular Exponentiation \hot}\label{sec:modexp}
\asked{Mid-sem 2024--25 Q5(ii): $100^{95} \bmod 9$, $658^{17} \bmod 17$, $145^{102} \bmod 101$. \new\ This year every RSA/ElGamal numerical also needs it (\S\ref{sec:rsa}).}

Modular exponentiation is exponentiation performed over a modulus: \imp{$a^b \bmod m$ means $(a^b) \bmod m$, the remainder of the whole power}; the modulus does not apply to the exponent. It is the core operation of public-key cryptography. RSA encrypts with $C = P^e \bmod n$ and decrypts with $P = C^d \bmod n$, Diffie--Hellman computes $g^x \bmod p$, ElGamal and DSS are built from the same operation, and elliptic-curve cryptography uses its additive analogue, double-and-add. The lecture introduces it with $2^{33} \bmod 30$ and $3^{100} \bmod 29$.

Computing the power first and reducing at the end fails for two separate reasons. The first is \textbf{size}: $3^{100}$ already has 48 digits, and with a 2048-bit exponent the power would have roughly $10^{616}$ digits, far beyond any memory. The second is \textbf{time}: multiplying by $a$ one step at a time needs $b - 1$ multiplications, which for $b \approx 2^{2048}$ is impossible. The size problem is solved by reducing modulo $m$ after every multiplication, which the product rule of \S\ref{sec:modular} allows and which keeps every intermediate value below $m^2$. The time problem is solved by repeated squaring, which needs only about $2\log_2 b$ multiplications. The lecture packages these ideas as two tools: reduce the base first, then build the power by repeated squaring, reducing after every multiplication.

\subsection{Tool 1: reduce the base first}
If $a \equiv r \pmod m$, then $a^b \equiv r^b \pmod m$. This follows from the product rule: a factor may be replaced by any congruent number, and $a^b$ is a product of $b$ factors equal to $a$. The base may therefore be replaced by any member of its residue class, and the best choice is the one with the smallest absolute value: the residue $r$ itself, or $r - m$ when $r > m/2$ (for example $23 \equiv -7 \pmod{30}$). Two special cases make some questions trivial. If $a \equiv 1$, then $a^b \equiv 1$ for every $b$. If $a \equiv -1$, then $a^b \equiv (-1)^b$, which is 1 for even $b$ and $-1 \equiv m - 1$ for odd $b$. A negative final result is turned into a residue by adding $m$.
\begin{center}\small
\begin{tabularx}{\textwidth}{@{}l l X l@{}}
\toprule
Problem & Base becomes & Working & Answer\\
\midrule
$23^3 \bmod 30$ & $23 \equiv -7$ & $(-7)^3 = 49 \times (-7) = -133 = -5 \times 30 + 17$ & \textbf{17}\\
$31^{500} \bmod 30$ & $31 \equiv 1$ & $1^{500} = 1$, however large the exponent & \textbf{1}\\
$242^{329} \bmod 243$ & $242 \equiv -1$ & $329$ is odd, so $(-1)^{329} = -1 \equiv -1 + 243$ & \textbf{242}\\
$11^{7} \bmod 13$ & $11 \equiv -2$ & $(-2)^7 = -128 = -10 \times 13 + 2$ & \textbf{2}\\
\bottomrule
\end{tabularx}
\end{center}
The first row can be checked by brute force, $23^3 = 12167 = 405 \times 30 + 17$; for large exponents no such check exists, which is why the method matters. The lecture reduces $-133$ in two steps ($-133 \to -13 \to -13 + 30 = 17$); formally $-133 = -5 \times 30 + 17$, and the result is the same.

\textbf{The exponent may not be reduced modulo $m$.} The exponent counts how many factors are multiplied; it is an ordinary integer, not an element of $\mathbb{Z}_m$, and the product rule says nothing about it. One counterexample settles it: $2^5 = 32 \equiv 2 \pmod 3$, but $5 \bmod 3 = 2$ and $2^2 = 4 \equiv 1 \pmod 3$. Powers do repeat, but with a different period: by Euler's theorem (\S\ref{sec:fermat}) $a^{\phi(m)} \equiv 1$ when $\gcd(a, m) = 1$, so the exponent may be reduced modulo $\phi(m)$ (modulo $p - 1$ for a prime $p$), never modulo $m$. Practice Q41 is built on this trap.

\subsection{Tool 2: repeated squaring}
When the base is not close to 0 or $\pm 1$, the first tool does not help: $88 \bmod 187$ is 88 or $-99$, and neither is small. The second tool attacks the number of multiplications instead. Write the exponent in binary, $b = \sum_i b_i 2^i$ with each $b_i \in \{0, 1\}$. Then
\[
a^b = \prod_{i \,:\, b_i = 1} a^{2^i}, \qquad\text{e.g.}\quad a^7 = a^4 \cdot a^2 \cdot a^1, \qquad a^{100} = a^{64} \cdot a^{32} \cdot a^4 ,
\]
and the powers $a^{2^i}$ are cheap because each is the square of the previous one: $a^{2^{i+1}} = \big(a^{2^i}\big)^2$. The lecture's procedure follows directly.
\begin{enumerate}[leftmargin=1.6em,itemsep=1pt]
  \item Write $b$ as a sum of powers of 2, which is its binary form: $7 = 4 + 2 + 1$, $100 = 64 + 32 + 4$.
  \item Build the table $a^1, a^2, a^4, \dots$ up to the largest power of 2 that does not exceed $b$. Each entry is the square of the previous one, reduced modulo $m$; use the negative representative whenever it is smaller in size.
  \item Multiply the entries whose exponents add up to $b$, reducing modulo $m$ after every multiplication.
\end{enumerate}
For $a^{100}$ this means 6 squarings ($a^2, a^4, \dots, a^{64}$) and 2 multiplications, 8 operations instead of 99. The general count is derived in the next subsection.

\subsubsection*{Example 1: $88^7 \bmod 187$ \quad ($7 = 4 + 2 + 1$)}
\begin{center}\small
\begin{tabular}{@{}lll@{}}
\toprule
Power & Computation & mod 187\\
\midrule
$88^1$ & & 88\\
$88^2$ & $88 \times 88 = 7744 = 41 \times 187 + 77$ & 77\\
$88^4$ & $77 \times 77 = 5929 = 31 \times 187 + 132$ & 132\\
\bottomrule
\end{tabular}
\end{center}
$88^7 = 88^4 \cdot 88^2 \cdot 88^1 \equiv 132 \cdot 77 \cdot 88$. The lecture multiplies all three first ($894{,}432 = 4783 \times 187 + 11$); reducing after each product keeps the numbers small: $132 \times 77 = 10164 \equiv 66$, then $66 \times 88 = 5808 \equiv \mathbf{11}$. The lecturer asks to bookmark this result. In the RSA lecture it is the encryption of $P = 88$ with $e = 7$, $n = 187$, and the homework $11^{23} \bmod 187 = 88$ (practice Q37) is the matching decryption with $d = 23$ (\S\ref{sec:rsa}).

\subsubsection*{Example 2: the last two digits of $29^5$}
The last $k$ digits of a positive integer $N$ are $N \bmod 10^k$, because $N = q \times 10^k + r$ where $r$ is the number formed by those digits. So the question asks for $29^5 \bmod 100$, with $5 = 4 + 1$:
\begin{center}\small
\begin{tabular}{@{}llll@{}}
\toprule
Power & Computation & mod 100 & or\\
\midrule
$29^1$ & & \textbf{29} & $-71$\\
$29^2$ & $29 \times 29 = 841$ & 41 & $-59$\\
$29^4$ & $41 \times 41 = 1681$ & 81 & $\mathbf{-19}$\\
\bottomrule
\end{tabular}
\end{center}
$29^5 = 29^4 \cdot 29^1 \equiv (-19)(29) = -551 \equiv -51 \equiv \mathbf{49}$. A second route confirms it: $29^2 \cdot 29^2 \cdot 29 \equiv 41 \cdot 41 \cdot 29 = 48749$, which ends in 49. In each row the representative with the smaller absolute value is used: 29 rather than $-71$, and $-19$ rather than 81.

\subsubsection*{Example 3: $3^{100} \bmod 29$ \quad ($100 = 64 + 32 + 4$, binary \ttx{1100100})}
\begin{center}\small
\begin{tabular}{@{}lllll@{}}
\toprule
Power & Computation & mod 29 & or & Use\\
\midrule
$3^1$ & & 3 & $-26$ & 3\\
$3^2$ & $3 \times 3$ & 9 & $-20$ & 9\\
$3^4$ & $9 \times 9 = 81$ & 23 & $-6$ & \textbf{$-6$} $\leftarrow$ needed\\
$3^8$ & $(-6)^2 = 36$ & 7 & $-22$ & 7\\
$3^{16}$ & $7 \times 7 = 49$ & 20 & $-9$ & $-9$\\
$3^{32}$ & $(-9)^2 = 81$ & 23 & $-6$ & \textbf{$-6$} $\leftarrow$ needed\\
$3^{64}$ & $(-6)^2 = 36$ & 7 & $-22$ & \textbf{7} $\leftarrow$ needed\\
\bottomrule
\end{tabular}
\end{center}
$3^{100} = 3^{64} \cdot 3^{32} \cdot 3^4 \equiv 7 \cdot (-6) \cdot (-6) = 252 = 8 \times 29 + 20$, so the answer is $\mathbf{20}$. The column of squares repeats ($-6 \to 7 \to -9 \to -6 \to 7$), because the powers of any element modulo $m$ eventually cycle. \emph{Preview of Fermat's theorem (\S\ref{sec:fermat}):} since 29 is prime, $3^{28} \equiv 1$, so $3^{100} = 3^{3 \cdot 28 + 16} \equiv 3^{16} \equiv 20$ \checkmark.

\subsubsection*{Example 4: $23^{16} \bmod 30$ (the exponent is a power of 2)}
When $b = 2^k$, the table is the whole computation: square $k$ times, reducing each time.
\[
23^{16} = \Big(\big((23^2)^2\big)^2\Big)^2:\quad 23 \equiv -7 \;\xrightarrow{\;^2\;}\; 49 \equiv 19 \equiv -11 \;\xrightarrow{\;^2\;}\; 121 \equiv 1 \;\xrightarrow{\;^2\;}\; 1 \;\xrightarrow{\;^2\;}\; \mathbf{1}
\]
Once a square reaches 1, every further square stays 1.

\textbf{Lecture homework:} (1) $11^{23} \bmod 187$; (2) the last three digits of $175^{209}$. Full solutions are in the practice section (Q37, Q38).

\subsection{Square-and-multiply: both scanning orders and the cost \hot}
\likely{``Compute $a^x \bmod n$ by square-and-multiply, showing every step'' (2--3 marks), inside every RSA/ElGamal numerical; sometimes ``write the algorithm and state its complexity''.}
The table method is how the computation is done by hand. Written as an algorithm it is called \textbf{square-and-multiply} (binary exponentiation), and it can scan the bits of the exponent in either direction. Both orders perform the same number of operations and give the same result; Forouzan prints the right-to-left version and Stallings the left-to-right one, so both are worth knowing. Below, $b = (b_{k-1} \dots b_1 b_0)_2$ has $k$ bits.
\begin{center}
\begin{minipage}[t]{0.47\textwidth}
\small\textbf{Right-to-left (Forouzan, Alg.\ 9.7)}\\[2pt]
\ttfamily\footnotesize
SqMul\_RL(a, b, n):\\
\hspace*{1em}y $\leftarrow$ 1\\
\hspace*{1em}for i $\leftarrow$ 0 to k$-$1:\\
\hspace*{2em}if b$_i$ = 1: y $\leftarrow$ (y $\times$ a) mod n\\
\hspace*{2em}a $\leftarrow$ (a $\times$ a) mod n\\
\hspace*{2em}\textrm{\emph{(not needed after the last bit)}}\\
\hspace*{1em}return y
\end{minipage}\hfill
\begin{minipage}[t]{0.47\textwidth}
\small\textbf{Left-to-right (Stallings)}\\[2pt]
\ttfamily\footnotesize
SqMul\_LR(a, b, n):\\
\hspace*{1em}f $\leftarrow$ 1\\
\hspace*{1em}for i $\leftarrow$ k$-$1 downto 0:\\
\hspace*{2em}f $\leftarrow$ (f $\times$ f) mod n\\
\hspace*{2em}if b$_i$ = 1: f $\leftarrow$ (f $\times$ a) mod n\\
\hspace*{1em}return f
\end{minipage}
\end{center}

\textbf{Right-to-left.} The bits are scanned from the least significant bit $b_0$ upwards, with two variables: $y$, the result so far (initially 1), and a running square, which holds $a^{2^i}$ at step $i$. If $b_i = 1$, the running square is multiplied into $y$; then it is squared for the next step. This is the lecture's table computed on the fly: the running square walks down the column $a, a^2, a^4, \dots$, and $y$ collects the rows that are needed. After step $i$ the invariant is $y = a^{(b \bmod 2^{i+1})}$, so after the last bit $y = a^b$. For $88^7 \bmod 187$ ($7 = 111_2$):
\begin{center}\small
\begin{tabular}{@{}ccll@{}}
\toprule
$i$ & $b_i$ & Multiplication (init.\ $y = 1$) & Squaring (init.\ $a = 88$)\\
\midrule
0 & 1 & $y = 1 \times 88 = 88$ & $a = 88^2 = 7744 \equiv 77$\\
1 & 1 & $y = 88 \times 77 = 6776 = 36 \times 187 + 44 \to 44$ & $a = 77^2 = 5929 \equiv 132$\\
2 & 1 & $y = 44 \times 132 = 5808 = 31 \times 187 + 11 \to \mathbf{11}$ & (not needed)\\
\bottomrule
\end{tabular}
\end{center}
The same answer as Example 1. Forouzan's own example, $17^{22} \bmod 21 = 4$, is traced in \S\ref{sec:primes}.

\textbf{Left-to-right.} The bits are scanned from the most significant bit down, with a single variable $f$: at every bit $f$ is squared, and if the bit is 1 it is also multiplied by $a$. To see why this works, let $c$ be the number formed by the bits processed so far (Stallings tracks it as an extra column). Initially $c = 0$ and $f = a^0 = 1$. Processing a bit appends it to $c$, which doubles $c$ and adds the bit: $c \leftarrow 2c + b_i$. Squaring $f$ doubles its exponent and multiplying by $a$ adds 1 to it, so the invariant $f \equiv a^c \pmod n$ holds after every step, and after the last bit $c = b$. For $3^{100} \bmod 29$ ($100 = 1100100_2$):
\begin{center}\small
\begin{tabular}{@{}cccll@{}}
\toprule
$i$ & $b_i$ & $c$ & Square & Multiply by 3 (if $b_i = 1$)\\
\midrule
6 & 1 & 1 & $1^2 = 1$ & $1 \times 3 = 3$\\
5 & 1 & 3 & $3^2 = 9$ & $9 \times 3 = 27$\\
4 & 0 & 6 & $27^2 = 729 = 25 \times 29 + 4 \to 4$ & --\\
3 & 0 & 12 & $4^2 = 16$ & --\\
2 & 1 & 25 & $16^2 = 256 = 8 \times 29 + 24 \to 24$ & $24 \times 3 = 72 \to 14$\\
1 & 0 & 50 & $14^2 = 196 = 6 \times 29 + 22 \to 22$ & --\\
0 & 0 & 100 & $22^2 = 484 = 16 \times 29 + 20 \to \mathbf{20}$ & --\\
\bottomrule
\end{tabular}
\end{center}
The result, 20, agrees with Example 3. The left-to-right order always multiplies by the same fixed base $a$, which is cheap when $a$ is small (primality tests often use $a = 2$) and needs only one variable. The right-to-left order computes the squarings independently of the bits, so they can be precomputed or run in parallel with the multiplications, and it matches the hand table.

\textbf{Cost: why it is $O(\log b)$.} Let the exponent have $k = \lfloor \log_2 b \rfloor + 1$ bits. Both orders perform one squaring per bit, of which one is free (squaring $f = 1$ at the start, or the unused last squaring), so $k - 1$ squarings in all. They perform one multiplication per 1-bit, at most $k$, and one of these is again free (a multiplication by 1). The total is therefore at most $2(k - 1) = 2\lfloor \log_2 b \rfloor$ modular multiplications, and about $1.5 \log_2 b$ on average, since about half the bits of a random exponent are 1. For $3^{100}$ (7 bits, three 1-bits) this is 6 squarings and 2 multiplications, 8 in all, against 99 for repeated multiplication. Each modular multiplication works on numbers below $n$ and needs one reduction, so with schoolbook arithmetic it costs $O((\log n)^2)$ bit operations, and the whole exponentiation costs $O(\log b \cdot (\log n)^2)$, which is polynomial in the length of the input. Repeated multiplication needs $b - 1 \approx 2^k$ steps, exponential in the length $k$ of the exponent. For a 2048-bit exponent the difference is about 4000 multiplications against $2^{2048}$.

This asymmetry is what public-key cryptography is built on: computing $a^x \bmod n$ is fast, while undoing it (the $e$-th root modulo $n$ in RSA, the discrete logarithm in Diffie--Hellman) has no comparably fast method (\S\ref{sec:primes}). Since the cost depends on the 1-bits, the common RSA public exponent $e = 65537 = 2^{16} + 1 = 10000000000000001_2$ needs only 16 squarings and one multiplication. The same dependence enables an attack: a multiplication happens only at a 1-bit, so the time or power consumed by a decryption leaks the bits of the private exponent $d$ unless the implementation runs in constant time (\S\ref{sec:rsa}).

\begin{mustknow}
$a^b \bmod m$: reduce the \textbf{base} modulo $m$ (to the smallest representative, possibly negative); \textbf{never reduce the exponent modulo $m$} (with $\gcd(a, m) = 1$ it may be reduced modulo $\phi(m)$, \S\ref{sec:fermat}).
Square-and-multiply: square at every bit, multiply at every 1-bit, reduce after every step; at most $2\lfloor \log_2 b \rfloor$ multiplications, i.e.\ $O(\log b)$.
\end{mustknow}

\begin{worked}[: 2024--25 Q5(ii), by base reduction and repeated squaring]
\textbf{Q.} Find (a) $100^{95} \bmod 9$ \quad (b) $658^{17} \bmod 17$ \quad (c) $145^{102} \bmod 101$.

\textbf{(a)} $100 = 11 \times 9 + 1$, so $100 \equiv 1 \pmod 9$ and $100^{95} \equiv 1^{95} = \mathbf{1}$.

\textbf{(b)} $658 = 38 \times 17 + 12$ and $12 \equiv -5 \pmod{17}$. The exponent is $17 = 16 + 1 = 10001_2$.
\begin{center}\small
\begin{tabular}{@{}lll@{}}
\toprule
Power & Computation & mod 17\\
\midrule
$12^1$ & & $12 \equiv -5$\\
$12^2$ & $(-5)^2 = 25 = 17 + 8$ & 8\\
$12^4$ & $8^2 = 64 = 3 \times 17 + 13$ & $13 \equiv -4$\\
$12^8$ & $(-4)^2 = 16$ & $16 \equiv -1$\\
$12^{16}$ & $(-1)^2 = 1$ & 1\\
\bottomrule
\end{tabular}
\end{center}
$658^{17} \equiv 12^{16} \cdot 12^1 \equiv 1 \times 12 = \mathbf{12}$.

\textbf{(c)} $145 = 101 + 44$, so the base is 44. The exponent is $102 = 64 + 32 + 4 + 2 = 1100110_2$.
\begin{center}\small
\begin{tabular}{@{}lll@{}}
\toprule
Power & Computation & mod 101\\
\midrule
$44^1$ & & 44\\
$44^2$ & $44^2 = 1936 = 19 \times 101 + 17$ & \textbf{17} $\leftarrow$\\
$44^4$ & $17^2 = 289 = 2 \times 101 + 87$ & $87 \equiv \mathbf{-14}$ $\leftarrow$\\
$44^8$ & $(-14)^2 = 196 = 101 + 95$ & $95 \equiv -6$\\
$44^{16}$ & $(-6)^2 = 36$ & 36\\
$44^{32}$ & $36^2 = 1296 = 12 \times 101 + 84$ & $84 \equiv \mathbf{-17}$ $\leftarrow$\\
$44^{64}$ & $(-17)^2 = 289 \equiv 87$ & $87 \equiv \mathbf{-14}$ $\leftarrow$\\
\bottomrule
\end{tabular}
\end{center}
Multiply the marked rows, reducing each time: $44^{64} \cdot 44^{32} \equiv (-14)(-17) = 238 = 2 \times 101 + 36 \to 36$; \;$36 \cdot 44^4 \equiv 36 \times (-14) = -504 = -5 \times 101 + 1 \to 1$; \;$1 \cdot 44^2 \equiv \mathbf{17}$.

\textbf{Answers: 1, 12, 17.} \textbf{Check} with Fermat's theorem (\S\ref{sec:fermat}): 17 and 101 are prime, so $12^{17} \equiv 12 \pmod{17}$, and $44^{100} \equiv 1 \pmod{101}$ gives $44^{102} \equiv 44^2 \equiv 17$ \checkmark.
\end{worked}

\begin{worked}[: square-and-multiply, left-to-right (Stallings' example)]
\textbf{Q.} Using the square-and-multiply algorithm, compute $7^{560} \bmod 561$. Show $c$ and $f$ after every step and count the multiplications.

\textbf{A.} $560 = 512 + 32 + 16 = 1000110000_2$ (10 bits). Start with $c = 0$, $f = 1$; at each bit square $f$, and multiply by 7 when $b_i = 1$.
\begin{center}\small
\begin{tabular}{@{}cccll@{}}
\toprule
$i$ & $b_i$ & $c$ & Square ($f^2 \bmod 561$) & $\times 7$ (if $b_i = 1$)\\
\midrule
9 & 1 & 1 & $1$ & $f = 7$\\
8 & 0 & 2 & $49$ & --\\
7 & 0 & 4 & $2401 = 4 \times 561 + 157 \to 157$ & --\\
6 & 0 & 8 & $24649 = 43 \times 561 + 526 \to 526$ & --\\
5 & 1 & 17 & $276676 = 493 \times 561 + 103 \to 103$ & $721 = 561 + 160 \to 160$\\
4 & 1 & 35 & $25600 = 45 \times 561 + 355 \to 355$ & $2485 = 4 \times 561 + 241 \to 241$\\
3 & 0 & 70 & $58081 = 103 \times 561 + 298 \to 298$ & --\\
2 & 0 & 140 & $88804 = 158 \times 561 + 166 \to 166$ & --\\
1 & 0 & 280 & $27556 = 49 \times 561 + 67 \to 67$ & --\\
0 & 0 & 560 & $4489 = 8 \times 561 + 1 \to \mathbf{1}$ & --\\
\bottomrule
\end{tabular}
\end{center}
\textbf{Answer: $7^{560} \bmod 561 = \mathbf{1}$}, using 9 non-trivial squarings and 2 non-trivial multiplications (11 instead of 559). The final $c = 560$ confirms the bits were read correctly. The result is notable: $561 = 3 \times 11 \times 17$ is composite, yet $7^{560} \equiv 1$ as if 561 were prime, because 561 is a Carmichael number, which fools the Fermat primality test (\S\ref{sec:primes}).
\end{worked}

\begin{trap}[: the usual slips]
\begin{itemize}[leftmargin=1.5em]
  \item Reducing the exponent modulo $m$, as in ``$117 \bmod 19 = 3$, so $5^{117} \equiv 5^3 \pmod{19}$''. Only the base is reduced modulo $m$.
  \item Reading the bits in the wrong direction. Right-to-left starts at the least significant bit; left-to-right starts at the most significant bit, and the $c$ column must end at $b$.
  \item Multiplying all the needed powers before reducing. It is not wrong, but the numbers grow quickly; reduce after every product.
  \item Writing $-7^3$ for $(-7)^3$. With an even exponent the brackets change the sign: $(-7)^2 = 49$ but $-7^2 = -49$. Give the final answer as a residue in $\{0, \dots, m-1\}$.
\end{itemize}
\end{trap}

\begin{fixbox}
\begin{itemize}
  \item All of Gemini's numbers are correct (checked by script: 17, 1, 242, 2, 11, 49, 20, 1). The two traces, Stallings' table for $7^{560} \bmod 561$ and the past-paper values 1, 12, 17 were also recomputed.
  \item Slide glitches noted by Gemini: in Example 2 the summary line briefly reads ``$88^7 \bmod 187 = 49$'' (copy--paste error; it should be $29^5 \bmod 100 = 49$). Example 3's slide stops at ``$252 \bmod 29$'' (the answer, 20, is only spoken).
  \item $-133 \bmod 30$: the lecture writes the intermediate step as $-13$ (truncated division), then adds 30. Formally $-133 = -5 \times 30 + 17$. Same answer.
  \item The slide writes $-7^3$ for $(-7)^3$. Strictly $-7^3 = -(7^3)$; the brackets matter when the exponent is even (e.g.\ $(-7)^2 = 49$ but $-7^2 = -49$).
\end{itemize}
\end{fixbox}
